
Художник Миккель Анжело садится на груду кирпичей и, подперев голову руками, начинает думать.

Вот проходит мимо петух и смотрит на художника Миккеля Анжело своими круглыми золотистыми глазами. Смотрит и не мигает.

Тут художник Миккель Анжело поднимает голову и видит петуха. Петух не отводит глаз, не мигает и не двигает хвостом.

Художник Миккель Анжело опускает глаза и замечает, что глаза что-то щиплет. Художник Миккель Анжело трет глаза руками. А петух не стоит уж больше, а уходит за сарай, на птичий двор к своим курам.

И художник Миккель Анжело поднимается с груды кирпичей, отряхивает со штанов красную кирпичную пыль, бросает в сторону ремешок и идет к своей жене.

По дороге художник Миккель Анжело встречает Комарова, хватает его за руку и кричит:

--- \textit{Смотри!}

Комаров смотрит и видит шар.

<<\textit{Что это?}>> --- шепчет Комаров.

А с неба грохочет: <<Это шар.>>

--- \textit{Какой такой шар?} --- шепчет Комаров.

А с неба грохот: <<\textit{Шар гладкоповерхностный!}>>

Комаров и художник Миккель Анжело садятся в траву и сидят, как грибы. Они держат друг друга за руки и смотрят на небо.

А на небе вырисовывается огромная ложка. Что же это такое? Никто этого не знает. Люди бегут и застревают в своих домах. Запирают двери и окна. Но разве это поможет? Куда там! Не поможет это.

Я помню, как в 1884--том году показалась на небе обыкновенная комета величиной с пароход. Очень было страшно. А тут ложка! Куда комете до такого явления.

Запирают окна и двери! 

Разве это может помочь? Против небесного явления доской не загородишься.

У нас в доме живет Николай Иванович Ступин, у него теория, что всё дым. А по-моему, не всё дым. Может, и дыма-то никакого нет. Ничего, может быть, нет. Есть одно только разделение. А может быть, и разделения-то никакого нет. Трудно сказать.

Говорят, один знаменитый художник рассматривал петуха и пришел к убеждению, что петуха не существует.

Художник сказал об этом своему приятелю, а приятель стал смеяться. Как же, говорит, не существует, когда он вот тут стоит, и я его отчетливо наблюдаю.

А великий художник опустил тогда голову и как стоял, так и сел на груду кирпичей.

\begin{center}
\textit{ВСЕ.}
\end{center}

\begin{flushright}
18 сентября 1934~г.
\end{flushright}

\begin{center}
* * *
\end{center}